% ----------------------
\subsection{Section 12 (May 1829)}
% ----------------------
	\label{DC12}
	
	% -----
	\subsubsection{Note on the JSP and BC versions}
	% -----
	
		The version featured in the JSP comes from BC, so for this section they are the same; the BC is therefore not included in the alternate versions below. This section doesn't appear at all in RB1 and is not listed in the index for that volume.
	
	% -----
	\subsubsection{Historical Background}
	% -----
		\label{DC12-HB}
		
		% --
		\paragraph{HDDC}
		% --
		
			``Joseph Smith and Oliver Cowdery received the Aaronic Priesthood from John the Baptist prior to the events that led to this revelation [Section 11] and also to Section 12. Therefore, these two sections should follow Section 13, which records their ordination'' (note 1, 218).
		
		% --
		\paragraph{JSP}
		% --
			
			``This revelation addressed Joseph Knight Sr., a Colesville, New York, resident and early supporter of JS. In late 1826, Knight began employing JS as a day laborer, and JS soon gained enough favor with Knight that he loaned JS his `horse and Cutter to go and see his girl,' Emma Hale, whom JS soon married. In September 1827, Knight was at the Joseph Smith Sr. home the night JS obtained the plates, and he later reported that JS gave him a detailed description of the plates and the Urim and Thummim. That winter JS and Emma, who had returned to Harmony, Pennsylvania, fell on hard times and asked Knight for assistance. `I let him have some little provisions and some few things out of the store a pair of shoes and three Dollars in money to help him a little,' wrote Knight.
			
			``Knight provided support that helped JS recommence the translation after Martin Harris lost the initial manuscript in July 1828. In January 1829, Joseph Smith Sr. and Samuel Smith stopped at Knight's home on their way from Manchester, New York, to visit JS in Harmony. Knight took them the rest of the way in his sleigh, and once they arrived in Harmony he gave Joseph Sr. a half-dollar and also gave JS money to buy paper for the translation. Knight may have also returned to Harmony to aid JS at other times between January and May 1829: Joseph Knight Jr. wrote, `Father and I often went to see him [JS] and carry him something to live upon.' In April 1829, the Knight family may have supplied all or part of the fifty-dollar payment that JS made on a thirteen-and-a-half-acre lot and small house purchased from his father-in-law, Isaac Hale.
			
			``When JS and Oliver Cowdery were running low on provisions in May 1829, they temporarily stopped translating and journeyed to Colesville to see Knight. Although absent at the time, Knight soon traveled to Harmony with more paper and with `provisions enough to Last till the translation was Done.' JS may have dictated this revelation during that visit. Introducing the revelation, JS's history stated that Knight was `very anxious to know his duty as to this work. I enquired of the Lord for him, and obtained as follows.'
			
			``This text addresses the same themes, largely in the same wording, as earlier revelations directed to Joseph Smith Sr. and Cowdery, alerting them to the emerging `great and marvelous work' and commanding them to `bring forth and establish the cause of Zion.' The revelation for Knight also broadened its potential audience by stating that it applied to `all those who have desires to bring forth and establish this work.' That the text itself does not specify Knight as its recipient also suggests that the call to work could have been to all who were willing to accept the responsibility.
			
			``The Book of Commandments contains the earliest copy of this revelation. John Whitmer did not include a copy in Revelation Book 1, indicating that JS may not have kept a copy or may have lost his copy. The editors of the Book of Commandments included the revelation, however, perhaps relying on a personal copy of the text kept by an early believer. When in May 1829 this revelation was dictated is unknown, in part because it is unclear how many times and when Knight traveled to Harmony to assist JS. The revelation dates from the same general period as the revelation for Hyrum Smith and may have preceded it. It is possible that Knight gave Hyrum a ride to Harmony, as he had done for Joseph Sr. and Samuel a few months earlier, and that the two revelations came during this trip. Neither Knight's history nor other accounts give exact dates as to when Knight arrived.''
			
	% -----
	\subsubsection{Versions}
	% -----
		\label{DC12-V}
	
		\begin{tabular}{p{1in}p{2in}p{1in}p{2in}}
			JSP		& \hyperref[EMS-1828-1829-JS-MAY-1829-B]{May 1829}	& BC & Chapter 11, 31 \\
			DC 1835	& Section 38, 169									& TS & 3:20 (15 August 1842), 884 \\
			MS		& 3:10 (February 1843), 164							& DC 1844 & Section 38, 250--251 \\
		\end{tabular}
		
		\medskip
		
		See the \href{https://drive.google.com/file/d/1goAvNz8jS4R8slYfMG_db55Ty2z_vmgK/view?usp=sharing}{sections comparison} document for more information about the version history.

	% -----
	\subsubsection{Section 12:1} % *DC12-1 
	% -----
		\label{DC12-1}

			\footnote{%
				% \label{}%
				BC has the following heading: ``1 \textit{A Revelation given to Joseph (K.,) [Joseph Knight Sr.] in Harmony, Pennsylvania, May, 1829, informing him how he must do, to be worthy to assist in the work of the Lord.}'' DC 1835 has, ``\textit{Revelation given to Joseph Knight, Sen. May, 1829.}'' DC 1844 has, ``\textit{Revelation given to Joseph Knight, Sen. May, 1829.}''
				}%
		A%
			\footnote{%
				% \label{}%
				More similarities here to \hyperref[DC6]{DC 6} and \hyperref[DC11]{DC 11}.
				}
		GREAT and marvelous work is about to come forth among the children of men.

		\begin{framed}
		\scriptsize
			\begin{compactdesc}
				\item[JSP] A GREAT and marvelous work is about to come forth among the children of men:
				\item[DC 1835] 1 A great and marvelous work is about to come forth among the children of men: 
				\item[DC 1844] 1 A great and marvellous work is about to come forth among the children of men: 
			\end{compactdesc}
		\end{framed}

	% -----
	\subsubsection{Section 12:2} % *DC12-2 
	% -----
		\label{DC12-2}

		Behold, I am God; give heed to my word, which is quick and powerful, sharper than a two-edged sword, to the dividing asunder of both joints and marrow; therefore, give heed unto my word.

		\begin{framed}
		\scriptsize
			\begin{compactdesc}
				\item[JSP] behold I am God, and give heed to my word, which is quick and powerful, sharper than a two-edged sword, to the dividing asunder of both joints and marrow: therefore, give heed unto my word.
				\item[DC 1835] behold I am God, and give heed to my word, which is quick and powerful, sharper than a two-edged sword, to the dividing asunder of both joints and marrow: therefore, given heed unto my word.
				\item[DC 1844] behold I am God, and give heed to my word, which is quick and powerful, sharper than a two-edged sword, to the dividing asunder of both joints and marrow: therefore, give heed unto my word.
			\end{compactdesc}
		\end{framed}

	% -----
	\subsubsection{Section 12:3} % *DC12-3 
	% -----
		\label{DC12-3}

		Behold, the field is white already to harvest;%
			\footnote{%
				% \label{}%
				Now picking up similarities to \hyperref[DC4]{DC 4}.
				}
		therefore, whoso desireth to reap let him thrust in his sickle with his might, and reap while the \hyperref[DAY]{day} lasts, that he may treasure up for his soul everlasting salvation in the kingdom of God.

		\begin{framed}
		\scriptsize
			\begin{compactdesc}
				\item[JSP] 2 Behold the field is white already to harvest, therefore whoso desireth to reap, let him thrust in his sickle with his might, and reap while the day lasts, that he may treasure up for his soul everlasting salvation in the kingdom of God:
				\item[DC 1835] 2 Behold the field is white already to harvest, therefore whoso desireth to reap, let him thrust in his sickle with his might, and reap while the day lasts, that he may treasure up for his soul everlasting salvation in the kingdom of God: 
				\item[DC 1844] 2 Behold the field is white already to harvest, therefore, whoso desireth to reap, let him thrust in his sickle with his might, and reap while the day lasts, that he may treasure up for his soul, everlasting salvation in the kingdom of God: 
			\end{compactdesc}
		\end{framed}

	% -----
	\subsubsection{Section 12:4} % *DC12-4 
	% -----
		\label{DC12-4}

		Yea, whosoever will thrust in his sickle and reap, the same is called of God.

		\begin{framed}
		\scriptsize
			\begin{compactdesc}
				\item[JSP] Yea, whosoever will thrust in his sickle and reap, the same is called of God: 
				\item[DC 1835] yea, whosoever will thrust in his sickle and reap, the same is called of God:
				\item[DC 1844] yea, whosoever will thrust in his sickle and reap, the same is called of God:
			\end{compactdesc}
		\end{framed}

	% -----
	\subsubsection{Section 12:5} % *DC12-5 
	% -----
		\label{DC12-5}

		Therefore, if you will ask of me you shall receive; if you will knock it shall be opened unto you.

		\begin{framed}
		\scriptsize
			\begin{compactdesc}
				\item[JSP] therefore, if you will ask of me you shall receive; if you will knock it shall be opened unto you.
				\item[DC 1835] therefore if you will ask of me you shall receive; if you will knock it shall be opened unto you.
				\item[DC 1844] ---therefore if you will ask of me you shall receive; if you will knock it shall be opened unto you.
			\end{compactdesc}
		\end{framed}

	% -----
	\subsubsection{Section 12:6} % *DC12-6 
	% -----
		\label{DC12-6}

		Now, as you have asked, behold, I say unto you, keep my commandments, and seek to bring forth and establish the cause of Zion.

		\begin{framed}
		\scriptsize
			\begin{compactdesc}
				\item[JSP] 3 Now as you have asked, behold I say unto you, keep my commandments, and seek to bring forth and establish the cause of Zion.
				\item[DC 1835] 3 Now as you have asked, behold I say unto you, keep my commandments, and seek to bring forth and establish the cause of Zion.
				\item[DC 1844] 3 Now as you have asked, behold I say unto you, keep my commandments, and seek to bring forth and establish the cause of Zion.
			\end{compactdesc}
		\end{framed}

	% -----
	\subsubsection{Section 12:7} % *DC12-7 
	% -----
		\label{DC12-7}

		Behold, I speak unto you, and also to all those who have desires to bring forth and establish this work;

		\begin{framed}
		\scriptsize
			\begin{compactdesc}
				\item[JSP] 4 Behold I speak unto you, and also to all those who have desires to bring forth and establish this work, 
				\item[DC 1835] 4 Behold I speak unto you, and also to all those who have desires to bring forth and establish this work,
				\item[DC 1844] 4 Behold I speak unto you, and also to all those who have desires to bring forth and establish this work, 
			\end{compactdesc}
		\end{framed}

	% -----
	\subsubsection{Section 12:8} % *DC12-8 
	% -----
		\label{DC12-8}

		And no one can assist in this work except he shall be humble and full of love,%
			\footnote{%
				% \label{}%
				These qualities appear in \hyperref[MOSIAH-3-19]{Mosiah 3:19} and \hyperref[ALMA-13-28]{Alma 13:28}.
				}
		having faith, hope, and charity, being temperate in all things,%
			\footnote{%
				% \label{}%
				See \hyperref[1CORINTHIANS-9-25]{1 Corinthians 9:25}.
				}
		whatsoever shall be entrusted to his care.

		\begin{framed}
		\scriptsize
			\begin{compactdesc}
				\item[JSP] and no one can assist in this work, except he shall be humble and full of love, having faith, hope and charity, being temperate in all things, whatsoever shall be intrusted to his care.
				\item[DC 1835] and no one can assist in this work, except he shall be humble and full of love, having faith, hope and charity, being temperate in all things, whatsoever shall be intrusted to his care.
				\item[DC 1844] and no one can assist in this work, except he shall be humble and full of love, having faith, hope and charity, being temperate in all things, whatsoever shall be intrus[ted] to his care.
			\end{compactdesc}
		\end{framed}

	% -----
	\subsubsection{Section 12:9} % *DC12-9 
	% -----
		\label{DC12-9}

		Behold, I am the light and the life of the world, that speak these words, therefore give heed with your \hyperref[MIGHT]{might}, and then you are called. Amen.

		\begin{framed}
		\scriptsize
			\begin{compactdesc}
				\item[JSP] 5 Behold I am the light and life of the world, that speaketh these words: therefore, give heed with your might, and then you are called: Amen.
				\item[DC 1835] 5 Behold I am the light and the life of the world, that speaketh these words: therefore, give heed with your might, and then you are called. Amen.
				\item[DC 1844] 5 Behold I am the light and the life of the world, that speaketh these words: therefore give heed with your might, and then you are called: Amen.
			\end{compactdesc}
		\end{framed}